% =========================================================
\section{Plan de trabajo y cierre}
% =========================================================

\begin{frame}{Cronograma}
  \begin{table}
  \fontsize{7.8}{9.0}\selectfont
  \begin{tabularx}{\textwidth}{>{\bfseries}l *{6}{>{\centering\arraybackslash}X}}
    \toprule
    Actividad & M1 & M2 & M3 & M4 & M5 & M6 \\
    \midrule
    Revisión y ajuste metodológico & \cellcolor{UCBlueSoft} & \cellcolor{UCBlueSoft} & & & & \\
    Desarrollo / experimentación & & \cellcolor{UCBlueSoft} & \cellcolor{UCBlueSoft} & \cellcolor{UCBlueSoft} & & \\
    Análisis de resultados & & & \cellcolor{UCYellow!35} & \cellcolor{UCYellow!35} & \cellcolor{UCYellow!35} & \\
    Escritura y defensa & & & & & \cellcolor{UCBlueSoft} & \cellcolor{UCBlueSoft} \\
    \bottomrule
  \end{tabularx}
  \end{table}
  \begin{UCKeyPoint}
    Reemplaza este cronograma por tu planificación real. Usa meses relativos o fechas concretas según lo que exija el seminario.
  \end{UCKeyPoint}
\end{frame}

\begin{frame}{Próximos pasos}
  \begin{columns}[T,onlytextwidth]
    \begin{column}{0.50\textwidth}
      \begin{UCBlock}{Trabajo técnico}
        \begin{itemize}
          \item Completar actividad pendiente 1.
          \item Completar actividad pendiente 2.
          \item Validar resultado o supuesto crítico.
        \end{itemize}
      \end{UCBlock}
    \end{column}
    \begin{column}{0.46\textwidth}
      \begin{UCBlock}{Trabajo de tesis de posgrado}
        \begin{itemize}
          \item Actualizar marco teórico.
          \item Consolidar resultados preliminares.
          \item Preparar capítulo o informe de avance.
        \end{itemize}
      \end{UCBlock}
    \end{column}
  \end{columns}
\end{frame}

\begin{frame}{Cierre}
  \begin{enumerate}
    \item \textbf{Contexto:} recordar el problema y su relevancia.
    \item \textbf{Avance:} resumir qué se completó y qué evidencia existe.
    \item \textbf{Limitaciones:} declarar supuestos, restricciones o incertidumbres.
    \item \textbf{Proyección:} indicar próximos pasos y decisiones esperadas.
  \end{enumerate}
  \begin{UCKeyPoint}
    Consejo: cerrar con tres o cuatro mensajes, cada uno conectado con una lámina anterior.
  \end{UCKeyPoint}
\end{frame}

\section{Referencias}

\begin{frame}[allowframebreaks]{Referencias}
  \nocite{*}
  \printbibliography
\end{frame}

\UCthanksframe[Gracias]

\appendix
\section{Anexos}

\begin{frame}{Anexo: respaldo técnico}
  \begin{itemize}
    \item Tabla de ejemplo con datos complementarios.
    \item Figura de ejemplo con validación o sensibilidad.
    \item Detalle metodológico que no cabe en la presentación principal.
    \item Referencias adicionales o supuestos utilizados.
  \end{itemize}
\end{frame}
